\section{Sharp ideal thresholds}
\label{sec:thresholds}

\subsection{Sufficiency}

\begin{proposition}[Summable side of the surface]
\label{prop:sufficiency}
If
\begin{equation}
  \sigma>1
  \quad\text{and}\quad
  \kappa_{b,q}b^{-\sigma}<1,
  \label{eq:sufficiency-conditions}
\end{equation}
then \(B_{b,s}\in S_q\).
\end{proposition}

\begin{proof}
By the two-sided unitary relation \cref{eq:two-sided-unitary}, it is enough
to take \(s=\sigma\).  Sum the embedded shell blocks first over a finite
square \(0\le k,\ell\le K\).  For the strict upper triangle, put
\(k=\ell+h\).  \Cref{eq:two-ratios} gives
\[
  \sum_{\ell\ge0}\sum_{h\ge1}
  \|B_{\ell+h,\ell}\|_{S_q}
  \le
  c_{b,q}
  \sum_{h\ge1}b^{h(1-\sigma)/2}
  \sum_{\ell\ge0}
  (\kappa_{b,q}b^{-\sigma})^\ell.
\]
Both geometric series converge under
\cref{eq:sufficiency-conditions}.  The lower triangle has the same block
norms.  On the diagonal,
\cref{eq:same-shell-factor,eq:weighted-comparison} gives
\[
  \sum_{k\ge0}\|B_{kk}\|_{S_q}
  \le
  \kappa_{b-2,q}
  \sum_{k\ge0}
  (\kappa_{b,q}b^{-\sigma})^k.
\]
Because \(q\ge1\), \(S_q\) is a Banach ideal.  The finite block sums are
therefore Cauchy in \(S_q\), and their limit has the prescribed matrix
entries.  This constructs the completed operator and proves the claim.
\end{proof}

\subsection{The universal wall}

\begin{proposition}[Column obstruction]
\label{prop:column-wall}
If \(\sigma\le1\), the formal matrix in \cref{eq:operator-intro} has no
bounded completion on \(\ell^2(\N)\).
\end{proposition}

\begin{proof}
The column indexed by \(n=1\) contains every \(m\) whose units digit lies
in \(\{0,\ldots,b-2\}\), because adding \(1\) creates no carry exactly on
that residue set.  The squared norm of this column is therefore at least
\[
  \sum_{\substack{m\ge1\\m\bmod b\in\{0,\ldots,b-2\}}}m^{-\sigma}.
\]
For \(0<\sigma\le1\), this is a finite union of arithmetic-progression
subseries of the divergent \(p\)-series; for \(\sigma\le0\), its summands
do not tend to zero.  Hence the column is not in \(\ell^2\), which is
incompatible with boundedness.
\end{proof}

The universal obstruction has no digit-spectrum input.  It remains the
active wall whenever \(\log_b\kappa_{b,q}\le1\); in that case the point
\(\sigma=\log_b\kappa_{b,q}\) is already on the unbounded side and is not a
second analytic endpoint.

\subsection{The higher-radix digit wall}

We use the standard orthogonal block pinching: if \((Q_j)\) are mutually
orthogonal projections, then
\[
  \Phi(T)=\bigoplus_j Q_jTQ_j
\]
is contractive on \(S_q\), \(1\le q<\infty\), and
\begin{equation}
  \|\Phi(T)\|_{S_q}^q
  =\sum_j\|Q_jTQ_j\|_{S_q}^q.
  \label{eq:pinching-norm}
\end{equation}
The statement follows first for finite families by averaging over diagonal
unitaries and then for countable families by monotone truncation.

Suppose \(\sigma>1\) and
\[
  \rho_{b,q}(\sigma):=\kappa_{b,q}b^{-\sigma}\ge1.
\]
For \(b\ge3\), pinch \(B_{b,\sigma}\) with the shell projections \(P_k\).
The lower half of \cref{eq:weighted-comparison} gives
\begin{equation}
  \|B_{kk}\|_{S_q}
  \ge
  b^{-\sigma}\kappa_{b-2,q}
  \rho_{b,q}(\sigma)^k.
  \label{eq:higher-radix-lower}
\end{equation}
Because \(\kappa_{b-2,q}>0\), the \(q\)th powers in
\cref{eq:higher-radix-lower} are not summable.  By
\cref{eq:pinching-norm}, \(B_{b,\sigma}\notin S_q\).

When \(\log_b\kappa_{b,q}>1\), this argument treats the active digit-wall
endpoint \(\sigma=\log_b\kappa_{b,q}\) and the entire interval strictly
below it but above the universal wall.  When
\(\log_b\kappa_{b,q}\le1\), no \(\sigma>1\) satisfies
\(\rho_{b,q}(\sigma)\ge1\), and \cref{prop:column-wall} already supplies
necessity.

\subsection{Binary paired-shell repair}
\label{sec:binary-pairing}

At \(b=2\), the preceding pinch is illegal as a digit-wall witness because
\(B_{kk}=0\) for every \(k\); see \cref{rem:binary-zero-block}.  We first
remove the phase by \cref{eq:two-sided-unitary} and work with the real
matrix \(B_{2,\sigma}\).  Let
\[
  \widehat P_j=P_{2j}+P_{2j+1},
  \qquad
  X_j:=B_{2j+1,2j}.
\]
The spaces \(I_{2j}\oplus I_{2j+1}\) are mutually orthogonal.  Since the
two diagonal-shell blocks vanish and \(B_{2,\sigma}\) is real symmetric,
\begin{equation}
  \widehat P_jB_{2,\sigma}\widehat P_j
  =
  \begin{pmatrix}
    0&X_j^*\\
    X_j&0
  \end{pmatrix}.
  \label{eq:binary-pair-block}
\end{equation}
Its nonzero singular values are those of \(X_j\), each repeated twice, so
\begin{equation}
  \|\widehat P_jB_{2,\sigma}\widehat P_j\|_{S_q}^q
  =2\|X_j\|_{S_q}^q.
  \label{eq:binary-double-sv}
\end{equation}

Using \cref{eq:binary-adjacent-exact} in the lower weighted comparison,
\begin{equation}
  \|X_j\|_{S_q}
  \ge
  2^{-3\sigma/2}
  (\kappa_{2,q}2^{-\sigma})^{2j}.
  \label{eq:binary-pair-lower}
\end{equation}
Thus, throughout the full binary digit-norm nonmembership range
\(\kappa_{2,q}2^{-\sigma}\ge1\) with \(\sigma>1\), the \(q\)th powers of
the paired compressions are nonsummable.  This includes both the
strict-below case, where the lower bound grows, and equality, where it
does not decay.  Equality is the exceptional repaired endpoint only
because same-shell pinching is unavailable.  As before, if the digit wall
is at or below \(1\), the universal column wall is the active boundary.

\begin{figure}[t]
  \centering
  \resizebox{\textwidth}{!}{\begin{tikzpicture}[
  >=Latex,
  panel/.style={
    draw=black!60,
    rounded corners=1.5mm,
    minimum width=67mm,
    minimum height=58mm,
    inner sep=3mm
  },
  shell/.style={
    draw=black!65,
    rounded corners=0.8mm,
    minimum width=12mm,
    minimum height=9mm,
    align=center,
    font=\scriptsize
  },
  active/.style={shell, fill=teal!14, very thick, draw=teal!65!black},
  zero/.style={shell, fill=gray!8, dashed},
  arr/.style={->, semithick, draw=purple!70!black},
]
  \node[panel] (leftpanel) at (-3.7,0) {};
  \node[panel] (rightpanel) at (3.7,0) {};

  \node[font=\small\bfseries, anchor=north] at (leftpanel.north)
    {Higher radices \(b\ge3\): same-shell pinch};
  \node[font=\small\bfseries, anchor=north] at (rightpanel.north)
    {Binary radix \(b=2\): paired-shell pinch};

  \node[active] (l0) at (-5.6,0.75) {\(I_k\)};
  \node[active, right=5mm of l0] (l1) {\(B_{kk}\)};
  \node[active, right=5mm of l1] (l2) {\(I_k\)};
  \draw[arr] (l0) -- (l1);
  \draw[arr] (l1) -- (l2);
  \node[align=center, font=\scriptsize, text width=59mm]
    at (-3.7,-1.18) {
      \(P_kBP_k=B_{kk}\), with\\[0.8mm]
      \(\|B_{kk}\|_{S_q}\gtrsim
      (\kappa_{b,q}b^{-\sigma})^k\).\\[0.8mm]
      Orthogonal shell projections force nonsummability
      when the ratio is at least one.
    };

  \node[zero] (r0) at (2.0,0.88) {\(I_{2j}\)};
  \node[zero, right=17mm of r0] (r1) {\(I_{2j+1}\)};
  \node[font=\scriptsize] at (r0.south) [below=1mm] {\(B_{2j,2j}=0\)};
  \node[font=\scriptsize] at (r1.south) [below=1mm] {\(B_{2j+1,2j+1}=0\)};
  \draw[arr, bend left=16] (r0) to node[above, font=\scriptsize] {\(X_j^*\)} (r1);
  \draw[arr, bend left=16] (r1) to node[below, font=\scriptsize] {\(X_j\)} (r0);
  \node[align=center, font=\scriptsize, text width=60mm]
    at (3.7,-1.28) {
      \(\widehat P_jB\widehat P_j=
      (\begin{smallmatrix}0&X_j\\X_j^*&0\end{smallmatrix})\),\quad
      \(\|\widehat P_jB\widehat P_j\|_{S_q}^q
      =2\|X_j\|_{S_q}^q\),\\[0.8mm]
      and \(\|X_j\|_{S_q}\gtrsim
      (\kappa_{2,q}2^{-\sigma})^{2j}\).
    };
\end{tikzpicture}
}
  \caption{The two digit-wall pinching mechanisms after phase removal.  For
  \(b\ge3\), nonzero same-shell blocks provide an orthogonal sequence.  For
  \(b=2\), \(X_j=B_{2j+1,2j}\) is the real weighted adjacent-shell block,
  and the mutually orthogonal two-shell compressions duplicate its
  singular values.  Both lower bounds reject ratio \(1\); the right-hand
  construction is required because every binary same-shell block is zero.}
  \label{fig:pinching-comparison}
\end{figure}

\subsection{Classification and operator classes}

\begin{theorem}[Sharp classification]
\label{thm:sharp-classification}
For every integer \(b\ge2\), every \(1\le q<\infty\), and every
\(s\in\C\), with \(\sigma=\Re s\),
\[
  B_{b,s}\in S_q
  \quad\Longleftrightarrow\quad
  \sigma>\max\{1,\log_b\kappa_{b,q}\}.
\]
\end{theorem}

\begin{proof}
The forward implication fails at \(\sigma\le1\) by
\cref{prop:column-wall}.  For \(\sigma>1\) but
\(\kappa_{b,q}b^{-\sigma}\ge1\), it fails by
\cref{eq:higher-radix-lower} when \(b\ge3\) and by
\cref{eq:binary-pair-lower} when \(b=2\).  Two-sided unitary equivalence
transfers these conclusions from real \(\sigma\) to complex \(s\).
The reverse implication is \cref{prop:sufficiency}.
\end{proof}

\begin{corollary}[Bounded, compact, Hilbert--Schmidt, and trace class]
\label{cor:operator-classes}
The following are equivalent:
\[
  B_{b,s}\text{ is bounded},\qquad
  B_{b,s}\text{ is compact},\qquad
  B_{b,s}\in S_2,\qquad
  \sigma>1.
\]
Furthermore,
\[
  B_{b,s}\in S_1
  \quad\Longleftrightarrow\quad
  \sigma>\alpha_b=\log_b\tau_b>1.
\]
\end{corollary}

\begin{proof}
For \(\sigma>1\), \cref{cor:digit-comparisons} gives
\(\log_b\kappa_{b,2}<1\), so \cref{thm:sharp-classification} yields
Hilbert--Schmidt, hence compact and bounded, membership.  For
\(\sigma\le1\), boundedness fails by \cref{prop:column-wall}.  The
trace-class statement is the case \(q=1\), together with
\(\tau_b>b\).
\end{proof}
